\documentclass[11pt]{article}
\usepackage[margin=1.1in]{geometry}
\usepackage{amsmath,amssymb,amsthm}
\usepackage[colorlinks=true,linkcolor=blue,citecolor=blue]{hyperref}

\newtheorem{theorem}{Theorem}[section]
\newtheorem{lemma}[theorem]{Lemma}
\newtheorem{proposition}[theorem]{Proposition}
\newtheorem{corollary}[theorem]{Corollary}
\newtheorem{definition}[theorem]{Definition}
\newtheorem{remark}[theorem]{Remark}
\newtheorem{example}[theorem]{Example}

\newcommand{\R}{\mathbb{R}}
\newcommand{\E}{\mathbb{E}}
\newcommand{\PP}{\mathbb{P}}
\newcommand{\emp}{\bar\mu_N}
\newcommand{\Sl}{S}
\newcommand{\vol}{\operatorname{vol}}
\newcommand{\diam}{\operatorname{diam}}

\title{The uniform doubling condition holds with high probability\\for independent samples}
\author{Yitzchak Shmalo\\
\small Einstein Institute of Mathematics, The Hebrew University of Jerusalem, Jerusalem, Israel\\
\small }
\date{June 2026}

\begin{document}
\maketitle

\begin{abstract}
The uniform doubling condition was introduced in \cite{ShmaloUDC} as a verifiable
property of a training set $T\subset\R^n$ guaranteeing stability of accuracy during
the training of deep neural networks with piecewise linear activation, refining the
doubling condition of Berlyand, Jabin and Safsten \cite{BJS}.  The stability theorems
of \cite{BJS,ShmaloUDC} are conditional: they apply to data satisfying the (uniform)
doubling condition, and the question of which data sets actually satisfy it was left
open --- indeed \cite{ShmaloUDC} notes that the condition cannot be fully verified
numerically even in small dimensions.  This paper proves that the condition is
satisfied \emph{with high probability} by independent samples.  Let $\mu$ be a
probability measure with density bounded above and below on a convex body
$K\subset\R^n$, and let $T$ consist of $N$ i.i.d.\ samples from $\mu$.  We show that
for every truncation number $k$ and every family $\mathcal L$ of truncation sets whose
truncation regions contain a ball of radius $r_0$, the training set $T$ satisfies the
uniform doubling condition with explicit constants
$(\kappa,\delta,\sigma,\beta)$ depending only on $n$, the density ratio, $r_0$ and
$\diam K$ --- not on $N$ --- and with scaling factor
$\ell_N=O\bigl((\log N/N)^{1/n}\bigr)$, with probability at least $1-N^{-1}$.  The
proof combines a new one-dimensional dichotomy for measures with log-concave profile
(every symmetric window, when widened by $\kappa$, either multiplies its mass by
$1+\sigma$ or already carries a fixed fraction of the total mass), the
Brunn--Minkowski inequality, a cone-volume estimate, and
Vapnik--Chervonenkis relative deviation bounds for classes of truncated slabs.  We
complement the theorem with two examples showing that the thickness hypothesis and
the scaling-factor decay rate are necessary in this generality.  As a corollary,
the stability-of-accuracy theorem of \cite{ShmaloUDC} applies with high probability
to i.i.d.\ training data.
\end{abstract}

\section{Introduction}\label{sec:intro}

A basic phenomenon in the training of deep neural networks (DNNs) with cross-entropy
loss is that the loss may decrease while classification accuracy deteriorates.
Berlyand, Jabin and Safsten \cite{BJS} identified a geometric condition on the
training data --- a \emph{doubling condition} on slabs in $\R^n$ --- which rules this
out: if the data satisfies the doubling condition, then a DNN with absolute-value
activation that achieves sufficiently high accuracy at some training time keeps high
accuracy for all later times.  In \cite{ShmaloUDC} the condition was made
\emph{uniform}: the constants were made independent of the choice of slabs, the class
of activations was broadened to piecewise linear activations such as Leaky ReLU, and
the stability theorems were proved under this uniform condition.  The condition reads
as follows (Definitions \ref{def:slab}--\ref{def:udc} below): for every slab
$S$ --- the region between two parallel hyperplanes, truncated by $k$ half-spaces
taken from a family $\mathcal L$ --- and for every width $w=\ell\kappa^i<\beta$ on the
geometric grid, the empirical mass of the widened slab $S_{\kappa w}$ must be at
least $\min\{\delta,\,(1+\sigma)\,\bar\mu(S_w)\}$.  Informally: no slab can trap a
small, isolated cluster of data; mass near a candidate decision boundary must keep
growing as the boundary is thickened, until a definite amount $\delta$ of data has
been captured.

The stability theorems of \cite{BJS,ShmaloUDC} are conditional on this property of
the data, and \cite{ShmaloUDC} observes that a full numerical verification is
infeasible already in small dimensions, resorting to a simplified surrogate.  This
leaves the basic question:

\begin{quote}
\emph{Which training sets actually satisfy the uniform doubling condition?}
\end{quote}

This paper answers the question for the most natural model of data: independent
samples from a distribution with density bounded above and below on a convex body.
We prove that such data satisfies the uniform doubling condition with high
probability, with constants independent of the sample size and with an explicit,
necessary decay of the scaling factor $\ell$.  Thus the hypothesis of the stability
theorems of \cite{BJS,ShmaloUDC} is not an idealization: it is a generic property of
i.i.d.\ data.

\subsection*{Main result}

Throughout, $K\subset\R^n$ is a convex body (compact, convex, nonempty interior),
$D:=\diam K$, and $\mu$ is a probability measure of the form
$d\mu=\rho\,dx$ on $K$ with
\begin{equation}\label{eq:density}
0<\rho_{\min}\le\rho(x)\le\rho_{\max}<\infty\quad(x\in K),
\qquad
\widetilde R:=\rho_{\max}/\rho_{\min}.
\end{equation}
The training set is $T=\{X_1,\dots,X_N\}$ with $X_i$ i.i.d.\ $\sim\mu$, and
$\emp:=\frac1N\sum_{i=1}^N\delta_{X_i}$ denotes the empirical measure, which is the
measure $\bar\mu$ of \cite{ShmaloUDC} normalized by $N$.

\begin{theorem}[Main theorem]\label{thm:main}
Let $\mu$ satisfy \eqref{eq:density}, let $k\in\mathbb N$, and let $r_0\in(0,D]$.
Set
\[
\kappa:=1+150\,\widetilde R,
\qquad
\sigma:=1,
\qquad
\delta:=\frac{3\,\rho_{\min}\,\omega_n r_0^{\,n}}{32\,\widetilde R},
\qquad
\beta:=D,
\]
where $\omega_n=\vol(B^n_1)$.  There are constants $C_0=C_0(n,k)$ and
$N_0=N_0(n,k,\widetilde R,r_0/D,\rho_{\min}r_0^n)$ such that for all $N\ge N_0$ the
following holds with probability at least $1-\frac1N$.  For \emph{every} family
$\mathcal L$ of truncation sets of size $k$ all of whose truncation regions
$G=H\cap K$ (Section~\ref{sec:prelim}) are either empty or contain a ball of radius
$r_0$, the training set $T$ satisfies the uniform doubling condition of
\cite[Definition 3.3]{ShmaloUDC} with doubling constant $\kappa$, mass cap $\delta$,
growth constant $\sigma$, scale parameter $\beta$, and global scaling factor
\[
\ell_N:=\Bigl(C_0\,\frac{\widetilde R}{\rho_{\min}}\Bigl(\frac{D}{r_0}\Bigr)^{n-1}
\frac{\log N}{N}\Bigr)^{1/n}.
\]
That is, for every truncated slab $S$ generated by a truncation set in $\mathcal L$
and every $i\in\mathbb N_0$ with $\ell_N\kappa^i<\beta$,
\[
\emp\bigl(S_{\ell_N\kappa^{i+1}}\bigr)\;\ge\;
\min\bigl\{\delta,\;(1+\sigma)\,\emp\bigl(S_{\ell_N\kappa^{i}}\bigr)\bigr\}.
\]
\end{theorem}

Two features deserve emphasis.  First, the constants $(\kappa,\delta,\sigma,\beta)$
do not depend on $N$: as the sample grows, the condition holds at finer and finer
scales ($\ell_N\to0$) with the \emph{same} doubling geometry, which is the regime
needed by the stability theorems.  Second, the family $\mathcal L$ may be infinite
--- the theorem is uniform over \emph{all} truncation sets with $r_0$-thick regions
--- so the same high-probability event serves every admissible network
simultaneously; in particular the truncation set of the DNN may move during training.

Both hypotheses are necessary in this generality, as the following examples
(Section~\ref{sec:sharp}) show.

\begin{example}[Thickness is necessary]\label{ex:sliver-intro}
For $K=[0,1]^n$ and the (population) measure $\mu=\mathrm{Leb}$, the truncation
region $H=\{x:x_2\le\varepsilon'\}$ produces, for the slab direction $e_2$, a
truncated slab whose mass saturates at $\varepsilon'$: for every
$\delta>\varepsilon'$ and every $\sigma,\ell>0$ the doubling condition fails at all
grid scales in $[\varepsilon',\beta)$.  Hence no choice of constants works uniformly
over all truncation sets, and a lower bound on the truncation-region thickness (or on
their masses) cannot be dispensed with.
\end{example}

\begin{example}[The scaling factor must decay no faster than a power]\label{ex:gap-intro}
For $K$ the Euclidean unit ball and $\mu$ uniform, with probability bounded below by
an absolute constant the sample contains a directional extreme point separated, in
that direction, from the rest of the sample by a gap of order $N^{-2/(n+1)}$.  A slab
of width below that gap captures this single point at all scales until the gap is
crossed, so the uniform doubling condition with $\sigma\le1$, $\delta\ge2/N$ forces
$\ell\ge c\,N^{-2/(n+1)}$.  Thus a sample-size--dependent scaling factor is
unavoidable, and the rate $\ell_N\asymp(\log N/N)^{1/n}$ of Theorem~\ref{thm:main}
is of the correct power-law type.
\end{example}

\subsection*{Method}

The proof has three independent ingredients.

\emph{(1) A one-dimensional dichotomy} (Lemma~\ref{lem:1d}): if $\nu$ is a finite
measure on $\R$ whose density is a log-concave function times a factor bounded
between $\rho_{\min}$ and $\rho_{\max}$, then for every center $c$ and width $w$,
widening the window $I_w(c)$ by the factor $\kappa\ge4$ either multiplies its mass by
$1+\frac{\kappa-1}{15\widetilde R}$ or the widened window already carries the mass
$\frac{3}{16\widetilde R}\nu(\R)$.  The proof is elementary but delicate; the
log-concavity enters through a quantitative ``no thin tails'' estimate.

\emph{(2) Geometry of truncated slabs.}  By Brunn--Minkowski, the slice profile of a
truncation region $G=H\cap K$ in any direction is log-concave, so ingredient (1)
applies to the projected measure and yields the doubling dichotomy for the
\emph{population} measure $\mu$, with the cap $\delta$ made uniform by the
$r_0$-thickness of $G$.  A cone-volume estimate (Lemma~\ref{lem:cone}) shows that any
truncated slab containing at least one sample point has, after widening, population
mass at least $c\,w^n$ --- the key to controlling scales where concentration fails.

\emph{(3) VC relative deviations.}  Truncated slabs are intersections of $k+2$
half-spaces, a class of VC dimension $O(nk\log k)$; the
Vapnik--Chervonenkis relative deviation bounds transfer the population dichotomy to
the empirical measure uniformly over all slabs and all grid scales at once, above the
scale $\ell_N$.  Below mass $\asymp\log N/N$ concentration is impossible, and it is
exactly here that the cone estimate and the choice of $\kappa$ absorb the
fluctuations: the population growth factor $1+\sigma_0=11$ beats the worst sampling
distortion, yielding the clean empirical constants $\sigma=1$ and an
$N$-independent mass cap $\delta$.

\subsection*{Relation to prior work}

The doubling condition originates in \cite{BJS}, where it is verified for measures
with density bounded above and below \emph{at the population level and for
non-truncated slabs with slab-dependent parameters}; the uniform, truncated version
that the stability theorems of \cite{ShmaloUDC} require is strictly stronger, and the
passage from population to empirical measure --- the content of this paper --- does
not appear there.  The empirical-process tools are classical
\cite{VC71,Vapnik,BEHW,AnthonyBartlett}; the geometric ingredient we need from convex
geometry is the Brunn--Minkowski inequality in its slice form
\cite{Gardner}.  For the role of the condition in the stability of accuracy, and for
the numerical surrogate it replaces, see \cite[Sections 3--5]{ShmaloUDC} and
\cite{BJS2}.

\subsection*{Organization}
Section~\ref{sec:prelim} recalls the definitions and the VC tools.
Section~\ref{sec:1d} proves the one-dimensional dichotomy.
Section~\ref{sec:population} proves the population doubling property.
Section~\ref{sec:empirical} proves Theorem~\ref{thm:main}.
Section~\ref{sec:sharp} contains the two sharpness examples, and
Section~\ref{sec:corollary} states the corollary for stability of accuracy.

\section{Preliminaries}\label{sec:prelim}

\subsection{Slabs, truncated slabs, and the uniform doubling condition}

We follow \cite[Section 3]{ShmaloUDC}.

\begin{definition}[Slab]\label{def:slab}
For a unit vector $u\in\R^n$, a center $s\in\R^n$ and a width $w>0$, the slab is
\[
S'_w(u,s):=\Bigl\{x\in\R^n:\ |(x-s)\cdot u|\le\frac w2\Bigr\}.
\]
\end{definition}

\begin{definition}[Truncated slab]\label{def:tslab}
Given additionally vectors $v_1,\dots,v_k\in\R^n$ and $t_1,\dots,t_k\in\R$, the
truncated slab is
\[
S_w=S_w(u,s,v_1,t_1,\dots,v_k,t_k):=S'_w(u,s)\cap H,
\qquad
H:=\{x:\ v_i\cdot x\le t_i,\ i=1,\dots,k\}.
\]
The half-space collection $h=\{(v_i,t_i)\}_{i\le k}$ is the \emph{truncation set};
a family of truncation sets of size $k$ is denoted $\mathcal L$.  We call
$G:=H\cap K$ the \emph{truncation region} of $h$ (relative to the data support $K$).
The width parameter $w$ scales the slab only; the truncation $H$ is fixed.
\end{definition}

\begin{definition}[{Uniform doubling condition, \cite[Definition 3.3]{ShmaloUDC}}]\label{def:udc}
Let $T\subset\R^n$ be a finite set with empirical measure $\emp$, and let
$\mathcal L$ be a family of truncation sets of size $k$.  The \emph{uniform doubling
condition} holds on $T$ with constants $\kappa>1$, $\delta,\sigma,\beta,\ell>0$ if
for every truncated slab $S$ generated by a truncation set $h\in\mathcal L$ and every
$i\in\mathbb N_0$ with $\ell\kappa^i<\beta$,
\[
\emp\bigl(S_{\ell\kappa^{i+1}}\bigr)\;\ge\;
\min\bigl\{\delta,\ (1+\sigma)\,\emp\bigl(S_{\ell\kappa^{i}}\bigr)\bigr\}.
\]
\end{definition}

In \cite{ShmaloUDC} the scaling factor may be chosen per slab; we prove the stronger
\emph{global} form, where the single scaling factor $\ell=\ell_N$ serves all slabs.

\begin{proposition}[Affine transfer of the uniform doubling condition]\label{prop:affine-transfer}
Let $T\subset\R^n$ satisfy the uniform doubling condition for a truncation family
$\mathcal L$ with constants $(\kappa,\delta,\sigma,\beta,\ell)$.  Let
$A(x)=Mx+b$ be an invertible affine map, and write
\[
\lambda=s_{\min}(M),\qquad \Lambda=s_{\max}(M).
\]
Let $A\mathcal L$ be the image truncation family, meaning that the half-space
$v\cdot x\le t$ is sent to
\[
(M^{-T}v)\cdot y\le t+(M^{-T}v)\cdot b .
\]
If
$\Lambda\ell<\lambda\beta/\kappa$, then $A(T)$ satisfies the uniform doubling
condition for $A\mathcal L$ with constants
\[
\bigl(\kappa^2,\delta,\sigma,\lambda\beta/\kappa,\Lambda\ell\bigr).
\]
If, in addition, the original doubling inequality holds at every width
$w\in[\ell,\beta)$ rather than only on the geometric grid, then the image
sample satisfies the same continuous-width form with constants
\[
\bigl(\kappa,\delta,\sigma,\lambda\beta,\Lambda\ell\bigr).
\]
\end{proposition}

\begin{proof}
We first interpolate the grid.  Suppose $w\in[\ell,\beta/\kappa)$.  Choose
$i\ge0$ with $\ell\kappa^i\le w<\ell\kappa^{i+1}$.  Monotonicity of slab mass
in the width and the grid doubling inequality give
\[
\emp(S_{\kappa^2 w})
\ge \emp(S_{\ell\kappa^{i+2}})
\ge \min\{\delta,(1+\sigma)\emp(S_{\ell\kappa^{i+1}})\}
\ge \min\{\delta,(1+\sigma)\emp(S_w)\}.
\]
Thus grid-UDC with multiplier $\kappa$ implies continuous-width UDC on
$[\ell,\beta/\kappa)$ with multiplier $\kappa^2$.

Now let $\bar\mu_A=A_\#\emp$ be the empirical measure of $A(T)$.  Take a slab
in the image with unit normal $u$, center $s$, and width $t$.  Its preimage is a
slab with unit normal
\[
v=\frac{M^Tu}{|M^Tu|}
\]
and width $w=t/|M^Tu|$, truncated by the corresponding member of $\mathcal L$.
Since
\[
\lambda\le |M^Tu|\le \Lambda,
\]
the range $t\in[\Lambda\ell,\lambda\beta/\kappa)$ implies
$w\in[\ell,\beta/\kappa)$.  Applying the interpolated inequality to the
preimage and using pushforward equality of empirical masses gives
\[
\bar\mu_A(S_{\kappa^2 t})
\ge \min\{\delta,(1+\sigma)\bar\mu_A(S_t)\}.
\]
This is exactly the UDC for $A(T)$ with multiplier $\kappa^2$, scale
$\Lambda\ell$, and upper parameter $\lambda\beta/\kappa$.  If the original
inequality is already available for every width in $[\ell,\beta)$, the same
argument applies directly with $t\in[\Lambda\ell,\lambda\beta)$ and multiplier
$\kappa$, giving the stated continuous-width constants.
\end{proof}

\subsection{VC tools}

For a class $\mathcal C$ of measurable subsets of $\R^n$, $\mathrm{vc}(\mathcal C)$
denotes its VC dimension.  Truncated slabs with $k$ truncations are intersections of
$k+2$ half-spaces:
\[
\mathcal C_{n,k}:=\Bigl\{\,\bigcap_{j=1}^{k+2}\{x:\ a_j\cdot x\le b_j\}\ :\
a_j\in\R^n,\ b_j\in\R\Bigr\}\supseteq\{\text{all truncated slabs}\}.
\]

\begin{lemma}[{\cite{BEHW}}]\label{lem:vc}
$\mathrm{vc}(\mathcal C_{n,k})\le 2(n+1)(k+2)\log_2\bigl(3(k+2)\bigr)=:d_{n,k}$.
\end{lemma}

\begin{lemma}[Relative uniform deviations; \cite{VC74,Vapnik,AnthonyBartlett}]\label{lem:relvc}
Let $\mathcal C$ be a class of measurable sets with $\mathrm{vc}(\mathcal C)=d\ge1$,
let $X_1,\dots,X_N$ be i.i.d.\ $\sim\mu$ with empirical measure $\emp$, and let
$\varepsilon\in(0,1)$.  Set
\[
u_N(\varepsilon):=\frac{4}{N}\Bigl(d\ln\frac{2eN}{d}+\ln\frac{8}{\varepsilon}\Bigr).
\]
Then with probability at least $1-\varepsilon$, simultaneously for all
$C\in\mathcal C$,
\begin{equation}\label{eq:relvc}
\mu(C)-\emp(C)\;\le\;\sqrt{\mu(C)\,u_N(\varepsilon)},
\qquad
\emp(C)-\mu(C)\;\le\;\sqrt{\emp(C)\,u_N(\varepsilon)} .
\end{equation}
\end{lemma}

\begin{proof}[Source]
These are the classical Vapnik--Chervonenkis relative deviation inequalities: each
one-sided supremum
$\sup_{C}\bigl(\mu(C)-\emp(C)\bigr)/\sqrt{\mu(C)}$ and
$\sup_{C}\bigl(\emp(C)-\mu(C)\bigr)/\sqrt{\emp(C)}$ exceeds $t$ with probability at
most $4\,\Pi_{\mathcal C}(2N)\,e^{-Nt^2/4}$, where
$\Pi_{\mathcal C}(2N)\le(2eN/d)^d$ by Sauer's lemma; see
\cite[Section A.5]{Vapnik}, \cite[Theorem 5.7]{AnthonyBartlett} and the survey
discussion in \cite{BBL}.  Solving
$4(2eN/d)^de^{-Nt^2/4}\le\varepsilon/2$ for $t^2$ and taking a union over the two
one-sided statements gives \eqref{eq:relvc} with the stated $u_N(\varepsilon)$.
\end{proof}

We record the two consequences used below.  On the event \eqref{eq:relvc}, for every
$C\in\mathcal C$ and every $\tau\ge u_N(\varepsilon)$:
\begin{align}
\mu(C)\ge16\tau
&\;\Longrightarrow\;
\emp(C)\ge\mu(C)-\sqrt{\mu(C)\,\tau}\ge\tfrac34\,\mu(C);
\label{eq:lower}\\
\emp(C)\ge4\tau
&\;\Longrightarrow\;
\mu(C)\ge\emp(C)-\sqrt{\emp(C)\,\tau}\ge\tfrac12\,\emp(C).
\label{eq:upper}
\end{align}
Indeed, in \eqref{eq:lower} $\sqrt{\mu\tau}\le\mu/4$; in \eqref{eq:upper}
$\sqrt{\emp\tau}\le\emp/2$.

\section{A one-dimensional dichotomy for log-concave profiles}\label{sec:1d}

For $c\in\R$ and $w>0$ write $I_w(c):=[c-\frac w2,\,c+\frac w2]$.

\begin{lemma}[Dichotomy]\label{lem:1d}
Let $\nu$ be a finite measure on $\R$ with density $g=\tilde\rho\,v$, where
$v:\R\to[0,\infty)$ is log-concave and integrable and
$\rho_{\min}\le\tilde\rho\le\rho_{\max}$ on $\{v>0\}$, with
$\widetilde R=\rho_{\max}/\rho_{\min}$.  Let $\kappa\ge4$.  Then for every $c\in\R$
and $w>0$,
\[
\nu\bigl(I_{\kappa w}(c)\bigr)\;\ge\;
\min\Bigl\{\frac{3}{16\widetilde R}\,\nu(\R),\
\Bigl(1+\frac{\kappa-1}{15\widetilde R}\Bigr)\,\nu\bigl(I_w(c)\bigr)\Bigr\}.
\]
\end{lemma}

\begin{proof}
Write $I=I_w(c)$, $\sigma_0=\frac{\kappa-1}{15\widetilde R}$.  We may assume
$\nu(I)>0$ and
\begin{equation}\label{eq:smallwindow}
\nu(I)<\frac{3}{16\widetilde R}\,\nu(\R),
\end{equation}
since otherwise $\nu(I_{\kappa w})\ge\nu(I)$ finishes the proof.  By
\eqref{eq:smallwindow} the complement of $I$ carries at least
$\bigl(1-\frac{3}{16}\bigr)\nu(\R)=\frac{13}{16}\nu(\R)$, so one of the two tails
carries at least $\frac{13}{32}\nu(\R)\ge\frac38\nu(\R)$.  All hypotheses are
invariant under reflection about $c$, so we may assume it is the right tail:
\begin{equation}\label{eq:Mplus}
M_+:=\nu\bigl(\bigl(c+\tfrac w2,\,\infty\bigr)\bigr)\;\ge\;\tfrac38\,\nu(\R)>0 .
\end{equation}
Put $p:=c+\frac w2$, $q:=c+\frac{\kappa w}2$, $E:=(p,q]$, and
$\Lambda:=q-p=\frac{(\kappa-1)w}2\ge\frac{3w}2$.  Since $\nu(I)>0$ and $M_+>0$, the
support interval of $v$ meets both $I$ and $(p,\infty)$; as it is an interval,
$v(p)>0$.  Let $\varphi:=-\log v$, a convex function on the support, and let
$m^*$ be a maximizer of $v$ (which exists: $v$ is log-concave, integrable, and not
a.e.\ zero; on the closure of its support $v$ is unimodal).  We will use twice the
following consequence of convexity: for $t_1<t_2<t_3$ in the support,
\begin{equation}\label{eq:slopes}
\frac{\varphi(t_2)-\varphi(t_1)}{t_2-t_1}\;\le\;
\frac{\varphi(t_3)-\varphi(t_2)}{t_3-t_2}.
\end{equation}

\medskip\noindent\emph{Case A: $m^*\le p$.}  Then $v$ is nonincreasing on
$[p,\infty)$.

\emph{A1: $v(q)\le v(p)/e$.}  Then $\varphi(q)-\varphi(p)\ge1$, so by
\eqref{eq:slopes} (with $t_1=p,t_2=q,t_3=t$), for every $t\ge q$ in the support,
$\varphi(t)\ge\varphi(q)+\frac{t-q}{\Lambda}$, i.e.\
$v(t)\le v(q)e^{-(t-q)/\Lambda}$.  Hence
\[
\int_q^\infty v\;\le\;\Lambda\,v(q)\;\le\;\int_p^q v,
\]
the last step because $v\ge v(q)$ on $[p,q]$.  Therefore
\[
M_+\;\le\;\rho_{\max}\Bigl(\int_p^qv+\int_q^\infty v\Bigr)
\;\le\;2\rho_{\max}\int_p^qv
\;\le\;2\widetilde R\,\nu(E),
\]
so by \eqref{eq:Mplus},
$\nu(I_{\kappa w})\ge\nu(E)\ge\frac{M_+}{2\widetilde R}\ge\frac{3}{16\widetilde R}\nu(\R)$:
the first branch of the dichotomy holds.

\emph{A2: $v(q)>v(p)/e$.}  Then $\varphi(q)-\varphi(p)<1$, so the slope of
$\varphi$ on $[p,q]$ is less than $1/\Lambda$; by \eqref{eq:slopes} (with
$t_2=p$, $t_3=q$ and $t_1=t$), every $t\le p$ in the support satisfies
$\varphi(p)-\varphi(t)\le\frac{p-t}{\Lambda}$, i.e.\
$v(t)\le v(p)\,e^{(p-t)/\Lambda}$.  For $t\in I$ we have $p-t\le w$ and
$w/\Lambda\le\frac2{\kappa-1}\le\frac23$, so
$\sup_Iv\le e^{2/3}v(p)$ and
\[
\nu(I)\;\le\;\rho_{\max}\,w\,e^{2/3}\,v(p).
\]
On the other hand $v\ge v(q)>v(p)/e$ on $E$, so
$\nu(E)\ge\rho_{\min}\Lambda\,v(p)/e$.  Hence
\[
\frac{\nu(E)}{\nu(I)}\;\ge\;\frac{\Lambda}{w}\cdot\frac{1}{\widetilde R\,e^{5/3}}
\;=\;\frac{\kappa-1}{2e^{5/3}\widetilde R}\;\ge\;\frac{\kappa-1}{11\widetilde R}
\;\ge\;\sigma_0,
\]
and $\nu(I_{\kappa w})\ge\nu(I)+\nu(E)\ge(1+\sigma_0)\nu(I)$: the second branch.

\medskip\noindent\emph{Case B: $m^*>p$.}  Then $v$ is nondecreasing on
$(-\infty,m^*]$; in particular $v\le v(p)$ on $I$, so
\begin{equation}\label{eq:IwB}
\nu(I)\;\le\;\rho_{\max}\,w\,v(p).
\end{equation}

\emph{B1: $m^*\ge q$.}  Then $v\ge v(p)$ on $E$, so
$\nu(E)\ge\rho_{\min}\Lambda v(p)$ and, by \eqref{eq:IwB},
$\nu(E)/\nu(I)\ge\frac{\kappa-1}{2\widetilde R}\ge\sigma_0$: second branch.

\emph{B2: $p<m^*<q$.}  Let $\lambda_1:=m^*-p$ and $\lambda_2:=q-m^*$, so
$\lambda_1+\lambda_2=\Lambda$.

If $\lambda_1\ge\Lambda/2$: on $[p,m^*]$ we have $v\ge v(p)$, so
$\nu(E)\ge\rho_{\min}\lambda_1v(p)\ge\rho_{\min}\frac\Lambda2v(p)$, and by
\eqref{eq:IwB}, $\nu(E)/\nu(I)\ge\frac{\kappa-1}{4\widetilde R}\ge\sigma_0$:
second branch.

If $\lambda_2\ge\Lambda/2$, consider the drop of $v$ over $[m^*,q]$, where $v$ is
nonincreasing.

\quad If $v(q)\le v(m^*)/e$: as in A1 (with $[m^*,q]$ in place of $[p,q]$),
$\int_q^\infty v\le\lambda_2v(q)\le\int_{m^*}^qv$, hence
\[
M_+\;\le\;\rho_{\max}\Bigl(\int_p^qv+\int_q^\infty v\Bigr)
\;\le\;2\rho_{\max}\int_p^qv\;\le\;2\widetilde R\,\nu(E),
\]
and as in A1 the first branch holds.

\quad If $v(q)>v(m^*)/e$: then on $[m^*,q]$, $v\ge v(q)>v(m^*)/e\ge v(p)/e$, so
$\nu(E)\ge\rho_{\min}\lambda_2v(p)/e\ge\rho_{\min}\frac{\Lambda}{2e}v(p)$, and by
\eqref{eq:IwB}, $\nu(E)/\nu(I)\ge\frac{\kappa-1}{4e\widetilde R}
\ge\frac{\kappa-1}{11\widetilde R}\ge\sigma_0$: second branch.

\medskip
All cases yield one of the two branches, completing the proof.
\end{proof}

\section{Population doubling for truncated slabs}\label{sec:population}

\begin{proposition}[Population dichotomy]\label{prop:population}
Let $\mu$ satisfy \eqref{eq:density}, let $\kappa\ge4$, and let $S_w$ be a truncated
slab (direction $u$, center $s$, truncation $H$) whose truncation region $G=H\cap K$
contains a ball of radius $r_0$.  Then for every $w>0$,
\[
\mu(S_{\kappa w})\;\ge\;\min\Bigl\{\delta_0,\
\Bigl(1+\frac{\kappa-1}{15\widetilde R}\Bigr)\mu(S_w)\Bigr\},
\qquad
\delta_0:=\frac{3\,\rho_{\min}\,\omega_nr_0^{\,n}}{16\,\widetilde R}.
\]
\end{proposition}

\begin{proof}
Let $\nu$ be the push-forward of $\mu|_G$ under $x\mapsto x\cdot u$; then
$\mu(S_w)=\nu\bigl(I_w(c)\bigr)$ with $c=s\cdot u$, for every $w$.  The measure
$\nu$ has density
\[
g(t)=\int_{G\cap\{x\cdot u=t\}}\rho\,d\mathcal H^{n-1}
\;\in\;\bigl[\rho_{\min}v(t),\ \rho_{\max}v(t)\bigr],
\qquad
v(t):=\vol_{n-1}\bigl(G\cap\{x\cdot u=t\}\bigr).
\]
$G$ is a convex body, so by the Brunn--Minkowski inequality
\cite[Section 5]{Gardner} the function $v^{1/(n-1)}$ is concave on the interval
$\{v>0\}$ when $n\ge2$; hence $v$ is log-concave (for $n=1$, $v$ is the indicator of
an interval, also log-concave).  Writing $g=\tilde\rho v$ with
$\tilde\rho:=g/v\in[\rho_{\min},\rho_{\max}]$ on $\{v>0\}$,
Lemma~\ref{lem:1d} applies and gives
\[
\mu(S_{\kappa w})\ge\min\Bigl\{\frac{3}{16\widetilde R}\,\mu(G),\
(1+\sigma_0)\,\mu(S_w)\Bigr\},
\]
and $\mu(G)\ge\rho_{\min}\vol\bigl(B(z,r_0)\bigr)=\rho_{\min}\omega_nr_0^n$ by the
thickness hypothesis.
\end{proof}

\begin{lemma}[Cone volume]\label{lem:cone}
Let $G\subseteq K$ be convex with $B(z,r_0)\subseteq G$, let $x\in G$, let $u$ be a
unit vector and $L>0$.  Then
\[
\vol\Bigl(G\cap\bigl\{y:\ |(y-x)\cdot u|\le L\bigr\}\Bigr)
\;\ge\;
b_n\Bigl(\frac{r_0}{D}\Bigr)^{n-1}\min(L,\,2r_0)^n,
\qquad
b_n:=\frac{\min\{\omega_n,\,\omega_{n-1}/n\}}{4^{\,n}},
\]
where $D=\diam K$ and, for $n=1$, the right-hand side is read as
$\tfrac12\min(L,2r_0)$.
\end{lemma}

\begin{proof}
For $n=1$: $G$ is an interval containing $x$ and an interval of length $2r_0$; the
segment from $x$ towards $z$ of length $\min(L,r_0)\ge\frac12\min(L,2r_0)$ lies in
$G\cap\{|y-x|\le L\}$.

Let $n\ge2$ and $\zeta:=|z-x|\le D$.

\emph{Case $\zeta\le L/2$.}  Every $y\in B(z,L/2)$ satisfies
$|(y-x)\cdot u|\le|y-z|+\zeta\le L$, so the slab contains $B(z,L/2)$ and
\[
\vol\bigl(G\cap\text{slab}\bigr)\ge\vol\bigl(B(z,\min(r_0,L/2))\bigr)
=\omega_n\min(r_0,L/2)^n\ge\omega_n4^{-n}\min(L,2r_0)^n,
\]
using $\min(r_0,L/2)\ge\frac14\min(L,2r_0)$ (check both cases $L\le2r_0$, $L>2r_0$).

\emph{Case $\zeta>L/2$.}  The convex hull $C$ of $\{x\}\cup B(z,r_0)$ lies in $G$ and
contains all points $x+s\hat a+e$ with $\hat a=(z-x)/\zeta$, $0\le s\le\zeta$ and
$|e|\le s\,r_0/\zeta$ (take $t=s/\zeta$ in $x+t(z-x)+tr_0b$, $|b|\le1$).  For
$0\le s\le L/2$ and $|e|\le s\,\min(1,r_0/\zeta)$ we get
$|(x+s\hat a+e-x)\cdot u|\le s+|e|\le2s\le L$, so this truncated cone lies in
$G\cap\text{slab}$ and
\[
\vol\ge\int_0^{L/2}\omega_{n-1}\bigl(s\,\min(1,r_0/\zeta)\bigr)^{n-1}ds
=\frac{\omega_{n-1}}{n}\,\min\Bigl(1,\frac{r_0}{\zeta}\Bigr)^{n-1}
\Bigl(\frac L2\Bigr)^n
\ge\frac{\omega_{n-1}}{n}\Bigl(\frac{r_0}{D}\Bigr)^{n-1}\Bigl(\frac L2\Bigr)^n,
\]
and $\bigl(\frac L2\bigr)^n\ge4^{-n}\min(L,2r_0)^n$ again in both cases
($L\le2r_0$: $(L/2)^n=2^{-n}\min^n$; $L>2r_0$: $(L/2)^n\ge r_0^n=2^{-n}\min^n$).
\end{proof}

\section{Proof of the main theorem}\label{sec:empirical}

\begin{proof}[Proof of Theorem~\ref{thm:main}]
Set $d:=d_{n,k}$ (Lemma~\ref{lem:vc}), $\varepsilon:=1/N$, and
\[
u_N:=\frac4N\Bigl(d\ln\frac{2eN}{d}+\ln(8N)\Bigr),
\qquad
c_*:=\rho_{\min}\,b_n\Bigl(\frac{r_0}{D}\Bigr)^{n-1},
\]
with $b_n$ from Lemma~\ref{lem:cone}.  Let $\mathcal E$ be the event of
Lemma~\ref{lem:relvc} for the class $\mathcal C_{n,k}$, so
$\PP(\mathcal E)\ge1-\frac1N$, and \eqref{eq:lower}--\eqref{eq:upper} hold on
$\mathcal E$ with $\tau=u_N$ for every truncated slab simultaneously --- note that
$\mathcal E$ does not depend on the family $\mathcal L$.  Define
\[
\ell_N:=\frac{2}{\kappa-1}\Bigl(\frac{32\,u_N}{c_*}\Bigr)^{1/n},
\]
which is bounded by the expression in the statement of the theorem for a suitable
$C_0(n,k)$, and let
\[
N_0:=\min\Bigl\{N:\ 32u_N\le\min\bigl\{c_*(2r_0)^n,\ 2\delta_0\bigr\}
\ \text{ and }\ \ell_N<\beta\Bigr\},
\]
with $\delta_0$ from Proposition~\ref{prop:population}; thus $N_0$ depends only on
$(n,k,\widetilde R,r_0/D,\rho_{\min}r_0^n)$.

Work on $\mathcal E$ and let $N\ge N_0$.  Fix a truncated slab $S$ generated by a
truncation set in $\mathcal L$, with truncation region $G$, and fix any scale
$w\in[\ell_N,\beta)$; this covers in particular every grid scale
$\ell\kappa^i<\beta$ for every $\ell\ge\ell_N$, so the theorem follows for the
scaling factor in the statement, which dominates $\ell_N$ as defined here.  We must
show
\begin{equation}\label{eq:goal}
\emp(S_{\kappa w})\;\ge\;\min\bigl\{\delta,\ 2\,\emp(S_w)\bigr\},
\qquad
\delta=\tfrac12\delta_0 .
\end{equation}

\emph{Case 1: $\emp(S_w)=0$.}  Then the right-hand side of \eqref{eq:goal} is $0$
and there is nothing to prove.

\emph{Case 2: $\emp(S_w)>0$.}  Then there is a sample point $x\in S_w\cap T$.
Almost surely $x\in K$, and $x\in H$ since $S_w\subseteq H$; hence $x\in G$, so $G$
is nonempty and, by hypothesis, contains a ball of radius $r_0$.  Since
$|(x-s)\cdot u|\le w/2$, the slab of half-width $L:=\frac{(\kappa-1)w}2$ around $x$
satisfies $\{y:|(y-x)\cdot u|\le L\}\subseteq S'_{\kappa w}(u,s)$, so by
Lemma~\ref{lem:cone},
\[
\mu(S_{\kappa w})\;\ge\;\rho_{\min}\vol\Bigl(G\cap\{|(y-x)\cdot u|\le L\}\Bigr)
\;\ge\;c_*\min(L,2r_0)^n .
\]
If $L\le2r_0$ then $\min(L,2r_0)^n=L^n\ge\bigl(\frac{(\kappa-1)\ell_N}2\bigr)^n
=\frac{32u_N}{c_*}$ by the definition of $\ell_N$ and $w\ge\ell_N$; if $L>2r_0$ then
$c_*\min(L,2r_0)^n=c_*(2r_0)^n\ge32u_N$ by the definition of $N_0$.  In both cases
\begin{equation}\label{eq:mass-big}
\mu(S_{\kappa w})\;\ge\;32\,u_N,
\end{equation}
hence by \eqref{eq:lower},
\begin{equation}\label{eq:emp-big}
\emp(S_{\kappa w})\;\ge\;\tfrac34\,\mu(S_{\kappa w})\;\ge\;24\,u_N .
\end{equation}

\emph{Case 2a: $\emp(S_w)\le12u_N$.}  Then by \eqref{eq:emp-big},
$\emp(S_{\kappa w})\ge24u_N\ge2\,\emp(S_w)$, proving \eqref{eq:goal}.

\emph{Case 2b: $\emp(S_w)>12u_N$.}  By \eqref{eq:upper},
$\mu(S_w)\ge\frac12\emp(S_w)$.  Proposition~\ref{prop:population} gives
$\mu(S_{\kappa w})\ge\min\{\delta_0,(1+\sigma_0)\mu(S_w)\}$ with
$\sigma_0=\frac{\kappa-1}{15\widetilde R}=10$ for our choice
$\kappa=1+150\widetilde R$ (note $\kappa\ge151\ge4$).

If $\mu(S_{\kappa w})\ge\delta_0$: since $\delta_0\ge16u_N$ (definition of $N_0$),
\eqref{eq:lower} gives
$\emp(S_{\kappa w})\ge\frac34\delta_0\ge\frac12\delta_0=\delta$, proving
\eqref{eq:goal}.

If $\mu(S_{\kappa w})\ge(1+\sigma_0)\mu(S_w)$: combining with
\eqref{eq:mass-big}, \eqref{eq:lower} applies to $S_{\kappa w}$ and
\[
\emp(S_{\kappa w})\;\ge\;\tfrac34(1+\sigma_0)\,\mu(S_w)
\;\ge\;\tfrac38(1+\sigma_0)\,\emp(S_w)
\;=\;\tfrac38\cdot11\cdot\emp(S_w)\;\ge\;2\,\emp(S_w),
\]
proving \eqref{eq:goal}.

Thus on $\mathcal E$, for every admissible $\mathcal L$, every truncated slab and
every grid scale below $\beta$, the doubling inequality holds with
$(\kappa,\delta,\sigma,\beta,\ell_N)$ and $\sigma=1$,
$\delta=\frac12\delta_0=\frac{3\rho_{\min}\omega_nr_0^n}{32\widetilde R}$.
\end{proof}

\begin{remark}
The proof gives more than stated: the doubling inequality holds at \emph{every}
scale $w\in[\ell_N,\beta)$, not only on the geometric grid; and the same event
$\mathcal E$ serves all truncation families simultaneously, so the truncation set of
a DNN may change during training without leaving the scope of the theorem.  The
constants have not been optimized.
\end{remark}

\section{Sharpness}\label{sec:sharp}

\begin{proposition}[Thickness is necessary]\label{prop:sliver}
Let $K=[0,1]^n$, $\mu=\mathrm{Leb}|_K$, and let $\kappa>1$,
$\delta,\sigma,\beta,\ell>0$ with $\kappa\ell<\beta$ be given.  For
$0<\varepsilon'<\min\{\delta,\,\beta/\kappa\}$ consider the truncation set
$h=\{(e_2,\varepsilon')\}$ (padded by $k-1$ copies), with truncation region
$G=\{x\in K:\ x_2\le\varepsilon'\}$.  Then the \emph{population} measure $\mu$
violates the doubling inequality for the truncated slab with direction $u=e_2$ and
center $s=(\tfrac12,\tfrac{\varepsilon'}2,\tfrac12,\dots)$ at every scale
$w\in[\varepsilon',\beta/\kappa)$:
\[
\mu(S_{\kappa w})=\varepsilon'
\;<\;\min\{\delta,\ (1+\sigma)\varepsilon'\}
=\min\{\delta,\ (1+\sigma)\mu(S_w)\}.
\]
Consequently the empirical measure of an i.i.d.\ sample violates it as well for all
large $N$, almost surely.  Hence no choice of constants in
Definition~\ref{def:udc} holds uniformly over all truncation sets: a quantitative
lower bound on the truncation regions, such as the $r_0$-thickness of
Theorem~\ref{thm:main} (here the inradius of $G$ is $\varepsilon'/2\to0$), is
necessary.
\end{proposition}

\begin{proof}
For $w\ge\varepsilon'$ the truncated slab
$S_w=\{x\in K:\ |x_2-\tfrac{\varepsilon'}2|\le\tfrac w2,\ x_2\le\varepsilon'\}$
equals $G$, so $\mu(S_w)=\mu(S_{\kappa w})=\varepsilon'$, while
$\min\{\delta,(1+\sigma)\varepsilon'\}>\varepsilon'$ because
$\delta>\varepsilon'$ and $\sigma>0$.  Some required grid scale lies in
$[\varepsilon',\beta)$: if $\ell\ge\varepsilon'$, take $i=0$, noting
$\ell<\beta$; if $\ell<\varepsilon'$, take the smallest $i$ with
$\ell\kappa^i\ge\varepsilon'$, which satisfies
$\ell\kappa^i<\kappa\varepsilon'\le\beta$ since $\varepsilon'<\beta/\kappa$.  At
that scale the doubling inequality fails.  For the empirical
statement, $\emp(S_w)\to\mu(S_w)=\varepsilon'$ and
$\emp(S_{\kappa w})\to\varepsilon'$ almost surely, while
$\min\{\delta,(1+\sigma)\emp(S_w)\}\to\min\{\delta,(1+\sigma)\varepsilon'\}
>\varepsilon'$.
\end{proof}

\begin{proposition}[The scaling factor must decay]\label{prop:gap}
Let $K=B^n_1$ be the unit ball, $\mu$ the uniform measure, $n\ge1$, and let
$\kappa>1$, $\sigma>0$, $\beta>0$ be fixed.  There are constants
$c_2(n,\kappa),c_3(n,\kappa)>0$ such that for all large $N$, with probability at
least $c_3(n,\kappa)$, the sample $T$ admits a (trivially truncated) slab witnessing the failure
of the doubling inequality at every scale
$w\in\bigl[\ell,\ c_2(n,\kappa)N^{-2/(n+1)}\bigr)$ whenever $\delta\ge\frac2N$.
Consequently, on an event of probability bounded below, the uniform doubling
condition with constants $(\kappa,\delta,\sigma,\beta)$, $\delta\ge\frac2N$, forces
\[
\ell\;\ge\;c_2(n,\kappa)\,N^{-2/(n+1)} .
\]
\end{proposition}

\begin{proof}
Let $F(h):=\mu(\{x:\ x_1>1-h\})$ be the cap mass; there are constants
$c_4,c_5(n)>0$ with $c_4h^{(n+1)/2}\le F(h)\le c_5h^{(n+1)/2}$ for $h\in(0,\tfrac12]$
(integrate the marginal density, which is comparable to $(1-t)^{(n-1)/2}$ near
$t=1$).  Fix $\Lambda:=2\kappa$ and let $h_N:=N^{-2/(n+1)}$.  Consider the events
\[
A:=\{\text{exactly one sample point has }x_1>1-h_N\},
\qquad
B:=\{\text{no sample point has }x_1\in(1-\Lambda h_N,\,1-h_N]\}.
\]
With $p_1:=F(h_N)\le c_5/N$ and
$p_2:=F(\Lambda h_N)-F(h_N)\le c_5\Lambda^{(n+1)/2}/N$,
\[
\PP(A\cap B)=N\,p_1\,(1-p_1-p_2)^{N-1}
\;\ge\;N\,\frac{c_4}{N}\Bigl(1-\frac{c_5(1+\Lambda^{(n+1)/2})}{N}\Bigr)^{N-1}
\;\ge\;c_3(n)>0
\]
for all large $N$.  On $A\cap B$, let $x^*$ be the unique sample point with
$x_1^*>1-h_N$; every other sample point has $x_1\le1-\Lambda h_N$, hence is at
$e_1$-distance greater than $(\Lambda-1)h_N$ from $x^*$.  Consider the slab with
$u=e_1$, center $s$ with $s\cdot e_1=x^*_1$, and trivial truncations.  For every
$w$ with $\kappa w<2(\Lambda-1)h_N$, the slabs $S_w$ and $S_{\kappa w}$ contain
$x^*$ and no other sample point, so
$\emp(S_{\kappa w})=\frac1N<\min\{\delta,\ (1+\sigma)\frac1N\}$, violating the
doubling inequality at scale $w$.  The violating range is
$w<c_2h_N$ with $c_2:=\frac{2(\Lambda-1)}{\kappa}=\frac{2(2\kappa-1)}{\kappa}\ge2$.
If $\ell<c_2N^{-2/(n+1)}=c_2h_N$, the grid scale $w=\ell$ itself lies in the
violating range (and $\ell<\beta$ for large $N$), so the condition fails.
\end{proof}

\begin{remark}
Propositions \ref{prop:sliver} and \ref{prop:gap} delimit
Theorem~\ref{thm:main} from both sides: the thickness hypothesis cannot be removed
(even for the population measure), and the scaling factor must be allowed to decay
at a power rate in $N$ (even for the most regular data).  The gap between the
necessary rate $N^{-2/(n+1)}$ and the sufficient rate $(\log N/N)^{1/n}$ of
Theorem~\ref{thm:main} is an interesting open problem.
\end{remark}

\section{Stability of accuracy for i.i.d.\ data}\label{sec:corollary}

Combining Theorem~\ref{thm:main} with the stability theorems of \cite{ShmaloUDC}
(Theorem 3.6 there, and its Leaky ReLU variants) yields an unconditional statement
for i.i.d.\ data: with probability at least $1-\frac1N$ over the sample, every DNN
of the class considered in \cite{ShmaloUDC} whose truncation set has
$r_0$-thick regions relative to $K$, and which at some training time $t_0$
classifies the data with sufficiently high accuracy in the sense required there,
retains the corresponding accuracy lower bound for all $t>t_0$, with the loss of
accuracy quantified by the constants
$(\kappa,\delta,\sigma,\beta,\ell_N)$ of Theorem~\ref{thm:main}.  We refer to
\cite[Section 3]{ShmaloUDC} for the precise stability statements, which take the
constants of the doubling condition as inputs; the contribution of the present
paper is that for i.i.d.\ data these inputs hold with overwhelming probability with
$N$-independent $(\kappa,\delta,\sigma,\beta)$ and explicitly decaying $\ell_N$.

Two directions seem worth pursuing.  First, closing the gap of
Proposition~\ref{prop:gap} would identify the critical decay of the scaling factor;
the cap analysis suggests that boundary effects dominate and that
$N^{-\Theta(1/n)}$ is not optimal for strictly convex bodies.  Second, the
log-concave profile structure used in Lemma~\ref{lem:1d} holds for every
log-concave measure $\mu$, not only for densities bounded above and below on a
convex body; extending Theorem~\ref{thm:main} to unbounded log-concave data
(e.g.\ Gaussian mixtures with separation) requires replacing the cone estimate of
Lemma~\ref{lem:cone}, and we expect the main dichotomy to persist.

\subsection*{Acknowledgments}
The research presented in this paper was supported by the European Research Council
(ERC) under the European Union's Horizon Europe research and innovation programme
(grant agreement No.\ 101041711), by the Israel Science Foundation (grant number
2258/19) and by the Israel Science Foundation (ISF Grant 4101/25).

\subsection*{Declaration on generative AI}
During the preparation of this work the author used generative AI tools to
accelerate drafting and revision.  All mathematical claims, proofs, and
bibliographic information were subsequently reviewed and edited by the author, who
takes full responsibility for the content of the manuscript.

\begin{thebibliography}{9}

\bibitem{BJS}
L.~Berlyand, P.-E.~Jabin and C.~A.~Safsten,
\emph{Stability for the training of deep neural networks and other classifiers},
Math.\ Models Methods Appl.\ Sci.\ \textbf{31} (2021), no.~11, 2345--2390.

\bibitem{BJS2}
L.~Berlyand and P.-E.~Jabin,
\emph{Mathematics of Deep Learning: An Introduction},
De Gruyter, Berlin, 2023.

\bibitem{ShmaloUDC}
Y.~Shmalo,
\emph{Stability of accuracy for the training of DNNs via the uniform doubling
condition},
Ann.\ Math.\ Artif.\ Intell.\ (2024), DOI 10.1007/s10472-023-09919-1;
arXiv:2210.08415.

\bibitem{VC71}
V.~N.~Vapnik and A.~Ya.~Chervonenkis,
\emph{On the uniform convergence of relative frequencies of events to their
probabilities},
Theory Probab.\ Appl.\ \textbf{16} (1971), 264--280.

\bibitem{VC74}
V.~N.~Vapnik and A.~Ya.~Chervonenkis,
\emph{Theory of Pattern Recognition},
Nauka, Moscow, 1974 (in Russian).

\bibitem{Vapnik}
V.~N.~Vapnik,
\emph{Statistical Learning Theory},
Wiley, New York, 1998.

\bibitem{BEHW}
A.~Blumer, A.~Ehrenfeucht, D.~Haussler and M.~K.~Warmuth,
\emph{Learnability and the Vapnik--Chervonenkis dimension},
J.\ Assoc.\ Comput.\ Mach.\ \textbf{36} (1989), no.~4, 929--965.

\bibitem{AnthonyBartlett}
M.~Anthony and P.~L.~Bartlett,
\emph{Neural Network Learning: Theoretical Foundations},
Cambridge University Press, Cambridge, 1999.

\bibitem{BBL}
S.~Boucheron, O.~Bousquet and G.~Lugosi,
\emph{Theory of classification: a survey of some recent advances},
ESAIM Probab.\ Stat.\ \textbf{9} (2005), 323--375.

\bibitem{Gardner}
R.~J.~Gardner,
\emph{The Brunn--Minkowski inequality},
Bull.\ Amer.\ Math.\ Soc.\ (N.S.) \textbf{39} (2002), no.~3, 355--405.

\end{thebibliography}

\end{document}
